\begin{lemma}[The structural identity $R_\gamma L^{1/2} u = Du$]
For every $u \in \operatorname{Dom}(L^{1/2})$ in the sense of
\noderef{FracPowDomain},
\[
  R_\gamma\bigl(L^{1/2}u\bigr) = Du
  \qquad\text{in } L^2(\gamma_n;\mathbb{R}^n),
\]
with $R_\gamma$ as in \noderef{RieszTransform}, $L^{1/2}$ as in \noderef{FracPow}
and $D$ as in \noderef{GaussD}.  This is the identity that the endpoint argument
exploits: applying $R_\gamma$ to the bad part $L^{1/2}u$ of a decomposition
produces the ordinary gradient of the potential $u$.
\end{lemma}
\begin{proof}
Write $u_\alpha$ for the Hermite coefficients of $u$ relative to
\noderef{HermiteBasis}, and put $w = L^{1/2}u$.

\emph{Step 1: the coefficients of $w$.}
Since $u \in \operatorname{Dom}(L^{1/2})$, the formula of \noderef{FracPow} with
$s = 1/2$ applies and gives
\[
  w_0 = 0, \qquad w_\alpha = |\alpha|^{1/2}\,u_\alpha \quad (\alpha \neq 0).
\]

\emph{Step 2: applying the pseudoinverse to $w$.}
By \noderef{FracPowCoeff} with $s = -1/2$ the multiplied family for $w$ is $0$ at
$\alpha = 0$ and, for $\alpha \neq 0$,
\[
  |\alpha|^{-1/2}\, w_\alpha = |\alpha|^{-1/2}|\alpha|^{1/2}u_\alpha = u_\alpha ,
\]
using $|\alpha| \geq 1 > 0$ from \noderef{MultiIndexSum}.  This family is square
summable, being dominated by $(|u_\alpha|^2)_\alpha$ which sums to $\|u\|_2^2$ by
Parseval (\noderef{HermiteBasis}); so $w \in \operatorname{Dom}(L^{-1/2})$ by
\noderef{FracPowDomain} and, by \noderef{FracPow},
$L^{-1/2}\Pi_0 w$ has coefficients $0$ at $\alpha = 0$ and $u_\alpha$ at
$\alpha \neq 0$.

By \noderef{FracPowCoeff} and \noderef{FracPow} with $s = 0$, those are exactly
the coefficients of $\Pi_0 u$, and the coefficient map of \noderef{HermiteBasis}
is injective, so
\[
  L^{-1/2}\Pi_0\bigl(L^{1/2}u\bigr) = \Pi_0 u .
\]
By \noderef{ProjZeroCentering}, $\Pi_0 u = u - \mathfrak{m}(u)\mathbf{1}$ where
$\mathfrak{m}(u) = \int u \, d\gamma_n$.

\emph{Step 3: applying $D$.}
By \noderef{RieszTransform}, $R_\gamma w = D\bigl(L^{-1/2}\Pi_0 w\bigr)$, which by
Step 2 equals $D\bigl(u - \mathfrak{m}(u)\mathbf{1}\bigr)$.

Now $u \in W^{1,2}(\gamma_n)$ by \noderef{SqrtDomainSobolev}, and
$-\mathfrak{m}(u)\mathbf{1} \in W^{1,2}(\gamma_n)$ with
$D(-\mathfrak{m}(u)\mathbf{1}) = 0$ by part (3) of \noderef{GaussDLinear}.  Part
(1) of \noderef{GaussDLinear}, applied to the pair
$u$ and $-\mathfrak{m}(u)\mathbf{1}$, therefore gives
$u - \mathfrak{m}(u)\mathbf{1} \in W^{1,2}(\gamma_n)$ and
\[
  D\bigl(u - \mathfrak{m}(u)\mathbf{1}\bigr)
    = Du + D\bigl(-\mathfrak{m}(u)\mathbf{1}\bigr) = Du + 0 = Du .
\]
Combining with Step 3, $R_\gamma(L^{1/2}u) = Du$.
\end{proof}
